% Part: computability
% Chapter: computability-theory
% Section: coding-computations

\documentclass[../../../include/open-logic-section]{subfiles}

\begin{document}

\olfileid{cmp}{thy}{cod}
\olsection{संगणनांचे सांकेतिकरण}

संगणनाच्या प्रत्येक प्रतिमानात पुढील गोष्टी करता येतात:
\begin{enumerate}
\item संगणनक्षम फलनांच्या \emph{व्याख्यांचे} पद्धतशीर वर्णन करता येते.
  उदाहरणार्थ, ट्यूरिंग यंत्रांचे विनिर्देश, पुनरावर्ती व्याख्या किंवा
  कार्यक्रमण भाषेतील कार्यक्रम अशा व्याख्या देतात असे मानता येते.
\item एखाद्या व्याख्येने दिलेल्या फलनाचे ठरावीक आदानासाठी झालेले
  संगणन पूर्णपणे नोंदवता येते. उदाहरणार्थ, संगणनाच्या प्रत्येक पायरीवरील
  विन्यासांची क्रमिका (यंत्राची अवस्था आणि पट्टीवरील मजकूर) ट्यूरिंग
  यंत्राचे संगणन वर्णन करते.
\item संगणनाची एखादी कथित नोंद खरोखरच, दिलेल्या व्याख्येच्या
  संगणनक्षम फलनाचे दिलेल्या आदानासाठी झालेले संगणन कसे घडले याची नोंद
  आहे का, हे तपासता येते (त्या आदानावर फलन परिभाषित आहे, म्हणजे
  संगणन थांबते).
\item संगणनाच्या अशा पूर्ण नोंदीच्या वर्णनातून दिलेल्या आदानासाठी
  फलनाचे मूल्य काढता येते. उदाहरणार्थ, थांबणाऱ्या ट्यूरिंग यंत्राच्या
  संगणनाच्या अगदी शेवटच्या पायरीत पट्टीवर असलेला मजकूर हेच मूल्य असते.
\end{enumerate}

सांकेतिकरण वापरून संगणनक्षम फलनाच्या प्रत्येक वर्णनाला संख्यात्मक
\emph{निर्देशांक} देता येतो, आणि त्या निर्देशांकावरून सूचना संगणनक्षम
रीतीने पुन्हा मिळवता येतात. त्याचप्रमाणे, संगणनाची पूर्ण नोंदही एका
संख्येने संकेतबद्ध करता येते. मग ``निर्देशांक~$e$ असलेल्या फलनाच्या
आदान~$x$ वरील संगणनाच्या नोंदीचा संकेतांक~$s$ आहे'' हा परिणामी
अंकगणितीय संबंध आणि ``संकेतांक~$s$ असलेल्या संगणनक्रमिकेचे निर्गत''
हे फलन संगणनक्षम असतात; किंबहुना ती आदिम पुनरावर्ती असतात.

हे मूलभूत तथ्य फार सामर्थ्यवान आहे. त्याच्या आधारे संगणनक्षमतेविषयी
अनेक महत्त्वाचे आणि आश्चर्यकारक परिणाम सिद्ध करता येतात, तेही निवडलेल्या
संगणन-प्रतिमानावर अवलंबून न राहता.

\end{document}
